% id: 42-08
% book: 베이직쎈 미적분1 · 02 함수의 연속
% section: 학교 시험 기출로 실전 감각 UP!
% figure: none
% answer: 
% answer_source: 
% variant_level: 0 (원본 전사)
% 이 파일은 build.mjs 가 items.json 에서 만든다. 고칠 때는 items.json 을 고친다.
\smash{\raisebox{1.5mm}{\dmkichul{베이직쎈 미적분1 42쪽 08번}}}
자전거를 타고 $\pt{P}$ 지점을 출발하여 $\pt{A}$ 지점에서 멈춘 다음 다시 $\pt{A}$ 지점을 출발하여 $\pt{B}$ 지점에 도착하였다. $\pt{P}$ 지점에서 $\pt{A}$ 지점까지 갈 때와 $\pt{A}$ 지점에서 $\pt{B}$ 지점까지 갈 때의 자전거의 최고 속력이 모두 $30\mkern3mu \mathrm{km/h}$일 때, $\pt{P}$ 지점에서 $\pt{A}$ 지점을 거쳐 $\pt{B}$ 지점에 갈 때까지 자전거의 속력이 $10\mkern3mu \mathrm{km/h}$인 순간은 적어도 $n$번이다. 이때 $n$의 값은?
\dmchoicesauto{}{$1$}{$2$}{$3$}{$4$}{$5$}
